\documentclass[11pt,a4paper]{article}
\usepackage{curso}
\begin{document}
\portada{Ejercicios --- Lección}{5.1}{Por qué el hardware digital es ineficiente para la IA probabilística}
\section{Ejercicios}
\begin{enumerate}[label=\textbf{E\arabic*.}]
\item (B) Estime la intensidad aritmética de una actualización de grado 16 con acoplamientos fp16 y estados de 1 byte.
\item (B) Z1: 269\,568 pbits, 50 MHz, $<1$ W. Acote la energía por actualización y compárela con $\kB T\ln2$ a 300 K.
\item (B) ¿Por qué la GPU queda ociosa en Gibbs disperso?
\item (B) Enumere los tres costos digitales que un pbit evita.
\item (I) Si los acoplamientos vinieran de DRAM en lugar de SRAM, ¿cómo cambia la estimación de 70 pJ?
\item (I) Un generador produce 64 bits aleatorios con 10 operaciones de 1 pJ. ¿Energía por bit? Compare con 300 aJ por bit de un generador físico.
\item (A) Mida en su GPU el tiempo y la energía (NVML) de $10^8$ actualizaciones de Gibbs en un grafo de grado 16 y calcule J/actualización.
\end{enumerate}
\section{Soluciones}
\begin{solucion}
\textbf{E1.} $\sim32$ FLOPs (16 multiplicaciones-suma) sobre $16\cdot2+16\cdot1+$ escritura $\approx50$ bytes: $\sim0{,}6$ FLOP/byte.

\textbf{E2.} $\le7{,}4\times10^{-14}$ J; $\kB T\ln2=2{,}87\times10^{-21}$ J: $\sim2{,}6\times10^7$ veces.

\textbf{E3.} Está limitada por ancho de banda: hay $\sim1$ FLOP por byte y poco reuso.

\textbf{E4.} Movimiento de datos lejanos, fabricación de aleatoriedad, márgenes de ruido.

\textbf{E5.} $12\times640$ pJ $\approx7{,}7$ nJ: dos órdenes de magnitud más.

\textbf{E6.} $10/64\approx0{,}16$ pJ $=156$ fJ por bit, unas 500 veces 300 aJ.

\textbf{E7.} Resultado típico: decenas a cientos de pJ por actualización (sistema), según la GPU y el empaquetado de datos.
\end{solucion}
\end{document}
